\section{La naturaleza del problema de la exploración}

\subsection{Por qué importa la exploración}

En el aprendizaje por refuerzo, el agente debe encontrar un equilibrio entre la \textbf{explotación} (exploitation, elegir la mejor acción conocida) y la \textbf{exploración} (exploration, probar acciones que podrían ser mejores).

\begin{figure}[H]
\centering
\includegraphics[width=0.9\textwidth]{slides-images/slide-04.jpg}
\caption{Ilustración intuitiva de la dificultad de la exploración: algunas tareas son sencillas para una persona, pero casi imposibles para un RL que parte de cero}
\end{figure}

\begin{importantbox}{El dilema central de la exploración}
Si siempre elegimos la acción con el valor estimado más alto en ese momento (política voraz), quizá nunca descubramos la política verdaderamente óptima. Por ejemplo, en el bandido multibrazo, si al accionar un brazo por primera vez se obtiene por casualidad una recompensa positiva, la política voraz seguirá accionando ese mismo brazo, aunque otros brazos tengan una recompensa esperada mayor.
\end{importantbox}

\subsection{Por qué Montezuma's Revenge se volvió una referencia clásica de exploración}

La clase empieza ilustrando el problema con un juego de recompensas escasas como Montezuma's Revenge: obtener la llave da recompensa y abrir la puerta da recompensa, pero morir a manos de la skull no necesariamente produce de inmediato una retroalimentación negativa lo bastante clara, y los comportamientos decisivos para completar la tarea están muy separados entre sí. Una persona sabe qué hacer no porque vea la reward, sino porque entiende la semántica de los sprites y la estructura de la tarea.

\begin{figure}[H]
\centering
\includegraphics[width=0.9\textwidth]{slides-images/slide-05.jpg}
\caption{Montezuma's Revenge: las recompensas escasas, las cadenas largas de comportamiento y la falta de comprensión semántica dificultan la exploración}
\end{figure}

\subsection{Ponerse en el lugar del algoritmo}

El ponente lleva la analogía más lejos con el juego de cartas Mao: solo sabes que «romper las reglas se castiga», pero las reglas mismas solo pueden descubrirse por ensayo y error, y no siempre son intuitivas. Esto describe con precisión la dificultad esencial de las tareas de horizonte largo con recompensas escasas.

\begin{figure}[H]
\centering
\includegraphics[width=0.9\textwidth]{slides-images/slide-06.jpg}
\caption{Analogía de Mao: las reglas solo se descubren por ensayo y error; cuanto más larga es la tarea y más ocultas están las reglas, más difícil resulta explorar}
\end{figure}

\subsection{Exploración y explotación son en realidad el mismo problema}

La clase subraya que dos formulaciones habituales son, en esencia, equivalentes:
\begin{itemize}
    \item «¿Cómo descubrir políticas de alta recompensa a las que solo se llega mediante comportamientos complejos y de cadena larga?»
    \item «¿Cuándo conviene seguir haciendo lo mejor que ya se conoce y cuándo conviene probar un comportamiento nuevo?»
\end{itemize}
La primera pone el acento en el temporally extended behavior; la segunda, en el exploration-exploitation trade-off; pero ambas preguntan lo mismo: \textbf{¿estamos dispuestos a aceptar hoy una pequeña pérdida a cambio de un rendimiento posiblemente mayor mañana?}

\begin{figure}[H]
\centering
\includegraphics[width=0.9\textwidth]{slides-images/slide-07.jpg}
\caption{Restaurantes, colocación de anuncios y prospección de pozos petroleros: ejemplos cotidianos de exploración y explotación}
\end{figure}

\begin{warningbox}{La exploración no es difícil solo porque el espacio de acciones sea grande}
Lo verdaderamente complicado es que muchos comportamientos de exploración decisivos no muestran beneficio a nivel local e incluso reducen el rendimiento a corto plazo. Por eso, si la señal de entrenamiento solo atiende al desempeño inmediato, el agent tiende a verse empujado hacia políticas locales conservadoras pero subóptimas.
\end{warningbox}

\begin{figure}[H]
\centering
\includegraphics[width=0.9\textwidth]{slides-images/slide-08.jpg}
\caption{Del bandit a los MDP grandes y sin estructura: la dificultad teórica del problema de la exploración crece rápidamente}
\end{figure}

\subsection{La exploración en las aplicaciones de RL}

\begin{itemize}
    \item \textbf{Robótica}: explorar distintas estrategias de agarre y distintos patrones de movimiento
    \item \textbf{Modelos de lenguaje grandes}: en RLHF/RL for reasoning, explorar distintas rutas de razonamiento
    \item \textbf{Sistemas de recomendación}: explorar contenido que podría gustarle al usuario pero que todavía no se le ha mostrado
\end{itemize}

\subsection{Resumen de la sección}

La exploración es uno de los retos fundamentales de todo problema de RL. Sin una exploración eficaz, el agente se queda atrapado en políticas subóptimas.
